\section{Introduction}
\label{sec:introduction}

Probabilistic epistemic models are usually asked two different questions. One concerns degrees: how much support does an agent assign to a proposition? The other concerns categories: does the agent believe, know, reject, or remain undecided? Threshold approaches connect the two by turning sufficiently high probability into all-or-nothing belief. Four-valued semantics adds a second axis. Positive and negative support can coexist or both fail, so the categorical result is no longer binary but belongs to
\[
  \{\Tval,\Fval,\Bval,\Nval\}
  =\{(1,0),(0,1),(1,1),(0,0)\},
\]
following the Belnap--Dunn/FDE tradition \cite{Belnap1977,Dunn1976}.

The 4-PEL project combines these ideas in a finite-accessibility evidence semantics. Positive and negative extensions are measured separately by the same local set function. A Lockean threshold is applied to each mass, producing a four-valued belief state. A separate evidence-stable knowledge operator then checks whether the complete FDE profile is invariant across the accessible range. This yields an exact factorization:
\[
  \Know_i\varphi=
  \begin{cases}
    \Bel_i\varphi,&\Stable_i(\varphi),\\
    \Fval,&\text{otherwise.}
  \end{cases}
\]
The knowledge operator is therefore not a second probability threshold. It is a qualitative stability filter on top of threshold belief.

The lightweight interface and its stronger probability-certified subclass are distinguished in Section~\ref{sec:kernel}. The paper studies one-step conditionalization and, separately, a chosen affine interpolation of its support endpoints:

\begin{question}[Dynamic phase geometry]
Which four-valued belief and knowledge endpoints can admissible conditionalization connect, and what additional phase structure follows if their signed supports are interpolated affinely?
\end{question}

The machine-checked answer develops in several layers. First, admissible conditionalization can both destroy and restore the stability gate. More strongly, a fixed six-world construction scheme realizes every ordered pair of four-valued values at the level of \(\Know(\Bel p)\). Thus the categorical reachability graph is complete. That global permissiveness is paired with a sharp local invariant: a belief value changes exactly when at least one of its two threshold decisions changes.

The four categorical states can consequently be treated as the vertices of a Boolean square. The number of changed threshold coordinates defines a wall count
\[
  \dthr\in\{0,1,2\}.
\]
Zero walls means the same FDE state, one wall means an edge move, and two walls mean a diagonal move. The two diagonals are \(\Tval\leftrightarrow\Fval\) and \(\Nval\leftrightarrow\Bval\).

The geometric interpretation can then be reconnected to the underlying rational support masses. A changed threshold bit is exactly endpoint straddling of the Lockean boundary. For the explicit affine interpolation
\[
  \gamma(t)=x+t(y-x),\qquad t\in\mathbb{Q},
\]
straddling produces a rational crossing parameter \(t\in[0,1]\) at which \(\gamma(t)=\thr\). A two-wall transition therefore contains two threshold events, one for positive support and one for negative support. Each affine crossing time is unique. Their order is positive-first, simultaneous, or negative-first, and the simultaneous case is exactly passage through the point \((\thr,\thr)\) in the support plane.

If the two crossings are not simultaneous, the rational midpoint between them occupies a categorical state adjacent to both diagonal endpoints. The order of the wall events determines the intermediate state. For example,
\[
  \Nval\to\Bval
\]
becomes either
\[
  \Nval\to\Tval\to\Bval
\]
or
\[
  \Nval\to\Fval\to\Bval.
\]
The same classification holds for the other diagonal orientations. Here ``time'' means the interpolation parameter, not elapsed time or an inferred intermediate stage of a one-step update. Wall count counts differing endpoint bits, not crossings along an arbitrary path.

A separate verified layer constructs complete probability-integrity-certified models. Nonnegative normalized rational point weights generate local measures satisfying an explicit finite-probability integrity contract. Convex interpolation preserves those distributions and makes the mass of every fixed event affine in the interpolation parameter. The local paths can then be packaged into complete strong 4-PEL models while worlds, accessibility, atomic valuation, and threshold remain fixed.

For fixed events this yields a compact structural picture; applying it to the support of a varying modal formula requires an additional bridge:
\[
\boxed{
\begin{gathered}
\text{probability simplex}\longrightarrow\text{strong model path}\\
\longrightarrow\text{affine fixed-event masses}\longrightarrow\text{threshold sides}\\
\longrightarrow\text{wall crossings}\longrightarrow\text{crossing order}\\
\longrightarrow\text{FDE phase sequence}.
\end{gathered}}
\]

\paragraph{Reliability refinement.}
Section~\ref{sec:mpfg} studies a second information-loss map: six reliability-tagged evidence cells project to four coarse cells. Projection preserves support and mixing, but not the recoverability of reliability. Fine modal stability is equivalent to coarse stability plus local fibre rigidity. These elementary semantic results are not a six-valued logical calculus; they extend the paper's investigation of which structure a coarse observation forgets.

\paragraph{Methodological discipline.}
The paper distinguishes four levels throughout. Machine-checked theorems are reported as such. Finite witness constructions are not silently generalized beyond their quantified statement. Geometric and philosophical readings are identified as interpretations. Literature and novelty claims remain deliberately modest until a systematic comparison has been completed.

\paragraph{Current boundary.}
Model-valued convex paths are now verified for weight-generated strong endpoint models on a common semantic skeleton. This does not yet imply that every legacy conditionalization witness has such a strong representation, that every intermediate point is itself generated by conditionalization of the source model, or that support events defined by arbitrary probability-sensitive modal formulas remain fixed along the path. The rational affine crossing theorem also remains distinct from an abstract intermediate-value theorem over arbitrary continuous paths. These formula-level, update-level, and topological questions are the next mathematical boundaries.
